\documentclass[12pt]{article}
\usepackage{amsmath, amsthm, amsfonts, mathtools, xcolor, caption}

\numberwithin{equation}{section}

\usepackage[style=numeric,sorting=none,giveninits=true]{biblatex}
\addbibresource{bibliography.bib}

\usepackage{etoolbox}
\usepackage[margin=1in]{geometry}

\DeclareMathAlphabet{\mymathbb}{U}{BOONDOX-ds}{m}{n}

\usepackage{comment}

%%%%%%%%%%%%%%%%%%%%%%%%%%%%%%%%%%%%%%%%%
%       Theorem Environments            %
%%%%%%%%%%%%%%%%%%%%%%%%%%%%%%%%%%%%%%%%%

\newtheorem{theorem}{Theorem}[section]
\newtheorem{lemma}[theorem]{Lemma}
\newtheorem{corollary}[theorem]{Corollary}
\newtheorem{proposition}[theorem]{Proposition}
\theoremstyle{definition}
\newtheorem{remark}{Remark}
\newtheorem{definition}{Definition}
\newtheorem{notation}{Notation}
\newtheorem{example}{Example}
\newtheorem{conjecture}[theorem]{Conjecture}
\newtheorem{openproblem}[theorem]{Open Problem}

%%%%%%%%%%%%%%%%%%%%%%%%%%%%%%%%%%%%%%%%%
%               New Commands            %
%%%%%%%%%%%%%%%%%%%%%%%%%%%%%%%%%%%%%%%%%

\newcommand{\ev}[1]{( #1 )} 
\newcommand{\ew}[1]{\left( #1 \right)} 

\newcommand{\N}{\mathbb{N}}
\newcommand{\Z}{\mathbb{Z}}
\newcommand{\C}{\mathbb{C}}
\newcommand{\Q}{\mathbb{Q}}
\newcommand{\R}{\mathbb{R}}



%%%%%%%%%%%%%%%%%%%%%%%%%%%%%%%%%%%%%%%%%
%               Document                %
%%%%%%%%%%%%%%%%%%%%%%%%%%%%%%%%%%%%%%%%%

\let\underbrace\LaTeXunderbrace
\let\overbrace\LaTeXoverbrace


\usepackage{parskip}

\begin{document}

\section*{Challenge 3: Caccetta--H\"aggkvist Conjecture}
\refstepcounter{section}

%\begin{definition}[Out-degree]
%Let $D$ be a simple digraph with vertex set $V(D)$ and edge set $E(D)$. The \emph{out-degree} $d^+(v)$ of a vertex $v$ is the number of edges $(v,w)\in E(D)$, and the \emph{minimum out-degree} of $D$ is the number $\min_{v\in V\ev{D}} d^+(v)$. 
%\end{definition}
%
%\begin{definition}[Directed cycle]
%A \emph{directed cycle of length $\ell\geq 2$} in a simple digraph $D$ is a sequence $v_0,\dots,v_{\ell-1}$ of distinct vertices such that $(v_i,v_{i+1})\in E(D)$ for every $i$, where the indices are taken modulo $\ell$.
%\end{definition}

Given a simple digraph, if each vertex has an out-neighbour, then it contains a directed cycle \cite{bangjensen2009digraphs}, but not necessarily a short one. As the minimum out-degree increases, girth is expected to decrease. A natural construction is the digraph on $n$ cyclically ordered vertices, where each vertex is joined to the next $r$ vertices. In this example, the minimum out-degree is $r$ and the girth is $\lceil n/r\rceil$. The Caccetta--H\"aggkvist conjecture asserts that no digraph with minimum out-degree $r$ can do better. 
\begin{conjecture}[Caccetta--H\"aggkvist Conjecture \mbox{\cite{caccetta1978minimal}}]\label{conj:ch}
Let $r\in \mathbb{N}$ with $r\geq 1$. Every simple directed graph on $n$ vertices with minimum out-degree at least $r$ has a cycle of length at most $\lceil n/r\rceil$.
\end{conjecture}
For $r=1$ the statement is elementary. The conjecture was proved for $r=2$ by Caccetta and H\"aggkvist \cite{caccetta1978minimal}, for $r=3$ by Hamidoune \cite{hamidoune1987note}, and for $r=4$ and $r=5$ by Ho\`ang and Reed \cite{hoang1987note}. Shen proved it for every $r$ under the additional assumption $n\geq 2r^2-3r+1$ \cite{shen2000girth}. 
%When $r$ is a constant fraction of $n$, the conjecture predicts that minimum out-degree at least $n/3$ forces a directed cycle of length at most $3$; it is known that minimum out-degree at least $(3-\sqrt{7})\,n\approx 0.354\,n$ suffices \cite{shen1998directed}. 
%See \cite{west_cacchagg} for an overview.

%%%%%%%%%%%%%%%%%%%%%%%%%%%%%%%%%%%%%%%%%%%%%%%%%%
\printbibliography
%%%%%%%%%%%%%%%%%%%%%%%%%%%%%%%%%%%%%%%%%%%%%%%%%%
\end{document}
